% !TEX root = ../main.tex
% anything appart from comments in this file will appear in the Research Section of the text.
% do your main writing here.
% use \autocite for bibliography and \ref to cite internal figures and tables

% here's how to add a section to the contents page.


% Template question
% What sources have you identified as relevant to this project, and how and why has this influenced your practical work?

% Grading matrix:
% Assessment criteria:
% Research: Establish the normative range of music in a particular style or genre.
%
% Intended learning outcomes:
% Conduct relevant, accurate and comprehensive research to establish the normative range of music in a particular genre or relevant to the work of a particular director, company or studio
%
% Goal (distinction):
% Detailed research into the style in question, referencing a wide range of source material. Clear evidence of critical reflection and analysis that informs the composition.

\phantomsection
\subsection*{Sourced soundtracks}
\addcontentsline{toc}{subsection}{Sourced soundtracks}

In researching this task, I looked for recent crime dramas, keeping an eye out for Nordic productions\footnote{Guy makes a mention of the ``Nordic Noir'' genre of crime dramas being the originators of the ``modern'' sound, so it felt worth looking out for.}.

The sources I used most actively are:
\begin{itemize}
    \item \textit{Archive 81} \autocite{archive} (2022)
    \item \textit{The Crow Girl} \autocite{crow} (2025)
    \item \textit{The Jackal} \autocite{jackal} (2025)
    \item \textit{A Nearly Normal Family} \autocite{family} (2023)
    \item \textit{Severance} \autocite{severance} (2022)
    \item \textit{The Watchful Eye} \autocite{eye} (2022)
    \item \textit{The Åre Murders} \autocite{are} (2025)
\end{itemize}

\phantomsection
\subsubsection*{Guy vs the industry}
\addcontentsline{toc}{subsection}{Guy vs the industry}

The brief states that Guy ``[doesn't] want traditional orchestral
instruments'', presumably based on his impression (in 2022) that orchestral instruments aren't much used in the genre. But when researching, I found that most of the soundtracks actually \emph{do} feature fairly prominent orchestral instruments, primarily strings and some brass. Examples include ``The Jackal'' \autocite{jackal-jackal} [00:55--01:30], ``A Possible Solution'' \autocite{familiy-solution} [01:00--01:30], ``End Credits (Extended Version)'' \autocite{eye-end} [00:56--01:30].

To comply with the brief, I have avoided any overt orchestral textures.

\phantomsection
\subsubsection*{Guy's criminal record\footnote{For general research into Guy's writing style, refer to the section ``Researching Guy Michelmore'' in \hyperref[appendixC]{appendix C}. This section is oriented towards crime-sounding tracks.}}
\addcontentsline{toc}{subsection}{Guy's criminal record}

I didn't find much by Guy that fit the genre very well, but some tracks from \textit{Kitti Katz} \autocite{kk} (2023) contain elements that I could draw on.

``Into the Dark'' \autocite{kk-into} has some notable similarities to the assessment track: a low pad, a screechy/noisy synth, and some sort of plucked string instrument [00:27--00:49]. It moves into a four-on-the-floor beat with kick and claps [00:50], before settling into a slower groove with a ride and snare [01:09]. These overt electronic drum beats have played a part in how I've layered the drums in my continuation with electronic kick and snare samples [01:04, 01:39].

\phantomsection
\subsection*{Influence on the track}
\addcontentsline{toc}{subsection}{Influence on the track}

\phantomsection
\subsubsection*{Synths and sound design}
\label{sec:synths}
\addcontentsline{toc}{subsection}{Synths and sound design}

The task brief states that the continuation should ``keep with the dark electronic vibe''.

\textit{Severance}'s ``Main Titles'' \autocite{sev-main} feature a bass synth that adds extra tension through pitch bending [00:37--01:07]. The sound of the bass is quite mellow, however, sounding like somewhere between a sine and a filtered triangle or saw wave. Interestingly, it shares some sonic similarities with a typical EDM reese bass, which I based part of my own bass sound on.

On Guy's end, ``Kat Fright''\autocite{kk-fright} from \textit{Kitti Katz} contains a section with a menacing synth bass that also employs pitch modulation [01:14].

I wanted to combine those influences with a slightly more distorted bass sound, similar to what ``Night Vision'' \autocite{are-night} from \textit{The Åre Murders} does [00:12].

The result is the bass that in the second half of the track [01:19 onwards]. It's a fairly complex bass sound consisting of multiple layers, one of which is a reese-like bass. While not as prominent as in ``Main Titles'', there's also a few pitch bends happening [01:19, 01:46].

On a different bit of sound design, ``Sophia'' \autocite{crow-sophia} from \textit{The Crow Girl} uses multiple layers of short sounds (vocal chops, scrapes, breaths, percussion) over a plucky bass to create tension [01:01--01:25]. The fact that some of the sounds sound closer (breaths in particular) while others sound further away, and that the sounds continuously change their position in the stereo field through panning, helps create a sense of dread and being surrounded. ``Sophia'' is a much darker track than the one Guy started, but I've drawn inspiration from it in the section after the first main theme statement, with the percussive tonal loop in the background with different notes panned differently, and the percussion and noise loops that are layered over it [00:40--01:03].

\phantomsection
\subsubsection*{Drums}
\addcontentsline{toc}{subsection}{Drums}

Drums aren't featured very prominently in many of the references I found, but Guy lays down a pretty strong beat, so it felt appropriate to leave them in.

``End Credits (Extended Version)'' \autocite{eye-end} from \textit{The Watchful Eye} features a heavily-processed drum beat with what sounds like a hard (distorted) 808 kick and a gated snare layered with a ticking percussion layer to create a sense of urgency and tension [00:54--01:25].

\textit{Severance}'s ``Main Titles'' \autocite{sev-main} also features both the ticking clock [00:22] and a backbeat [00:37--01:07] with what sounds like an electronic kick (although the beat is much less processed than in ``End Credits''). The effect here is less one of urgency, but does emphasize passing time\footnote{A deeper analysis of what these elements represent in Severance is out of scope for this commentary.}.

In my track, the ticking clock first makes its appearance in the theme continuation [00:32]. The track has a faster pulse than both  ``End Credits (Extended Version)'' and ``Main Titles'', however, so the effect is less pronounced, and it's also not featured as prominently.

Regarding the grooves, Guy's base sounds fairly acoustic, consisting of a drum
loop sample and a sampled acoustic drum kit. I didn't want to go into overly
processed electronic drum samples, but instead opted for something in between
``End Credits'' and ``Main Titles''. The beat that underpins the repeat of the A
section is filtered and processed acoustic-based backbeat [01:04--01:31].
Because the start of this section is pretty sparse, this is also where the
ticking clock is the most pronounced. For the return to the main theme [01:40],
the drums return to the syncopated pattern established during the first full
theme statement [00:24], but to up the ante, I've layered in more hard-hitting
kicks and snares, similar to ``End Credits (Extended Version)''.

\phantomsection
\subsubsection*{Abrupt ending}
\addcontentsline{toc}{subsection}{Abrupt ending}

As heard in ``Archive 81 -- Titles'' \autocite{archive-titles} from \textit{Archive 81}, ``Main Titles'' \autocite{sev-main} from \textit{Severance}, and ``End Credits (Extended Version)'' \autocite{eye-end} from \textit{The Watchful Eye}, ending on a crescendo of noise that cuts off abruptly is a modern trope in this genre of music.

Based on input from René Boscio in a 121 session, I opted for ending my track in the same way [02:00--02:04].
